The adaptive-ball identity is definitional: mutual equivalence is the conjunction of its one-sided support condition and membership in \(\mathcal S_{\mu_X}(\boldsymbol P)\), with common conull sets replaced by their finite intersection.

For a mixture identity \(Q_a=e_aP_a+(1-e_a)R_a\), the inequality \(Q_a\ge e_aP_a\) and \(e_a\ge\kappa>0\) imply \(P_a\ll Q_a\).  Moreover,
\[
 Q_a(x,\cdot)\ll P_a(x,\cdot)\quad\Longleftrightarrow\quad R_a(x,\cdot)\ll P_a(x,\cdot). \tag{1}
\]
The reverse implication follows by addition.  For the forward implication, \((1-e_a(x))R_a(x,\cdot)\le Q_a(x,\cdot)\) and \(1-e_a(x)\ge\kappa>0\).  Intersecting the two armwise conull sets proves the mixture-class identity.

Finally, dropping reverse support from any mutual SCM witness gives a one-sided witness.  Conversely, take a one-sided witness whose \(\boldsymbol Q\) belongs to \(\mathcal S_{\mu_X}(\boldsymbol P)\).  Conditional one-sided completeness gives \(P_a\ll Q_a\) on a common conull set; together with \(Q_a\ll P_a\), the same witness satisfies the conditional reverse-support assumption.  This proves the SCM identity and the final automatic-support assertion.